\documentclass[12pt,oneside]{book}

% -------------------------------------------------
% PAQUETES
% -------------------------------------------------
\usepackage[spanish,es-nodecimaldot]{babel}
\usepackage[utf8]{inputenc}
\usepackage[T1]{fontenc}
\usepackage{lmodern}
\usepackage[a4paper,margin=2.5cm]{geometry}
\usepackage{microtype}
\usepackage{setspace}
\usepackage{graphicx}
\usepackage{xcolor}
\usepackage{fancyhdr}
\usepackage{titlesec}
\usepackage{hyperref}
\usepackage{bookmark}
\usepackage{tocloft}
\usepackage{tcolorbox}
\tcbuselibrary{theorems, skins, breakable}
\usepackage{amsmath,amssymb}
\usepackage{enumitem}
\usepackage{csquotes}

% -------------------------------------------------
% CONFIGURACIÓN VISUAL
% -------------------------------------------------
\definecolor{azulCosmos}{HTML}{1B3A6B}
\definecolor{azulCielo}{HTML}{2980B9}
\definecolor{doradoEstrella}{HTML}{D4A017}
\definecolor{grisClaro}{HTML}{F4F6F8}
\definecolor{grisTexto}{HTML}{2E2E2E}
\definecolor{grisSec}{HTML}{6C757D}
\definecolor{colorConcepto}{HTML}{EBF5FB}
\definecolor{colorActividad}{HTML}{EAFAF1}
\definecolor{colorIdea}{HTML}{FEF9E7}
\definecolor{colorNota}{HTML}{F2F3F4}
\definecolor{colorCompromiso}{HTML}{EAF0FB}

\hypersetup{
    colorlinks=true,
    linkcolor=azulCosmos,
    urlcolor=azulCielo,
    citecolor=azulCosmos,
    pdftitle={Clase de Astrobiología: Ecuación de Drake},
    pdfauthor={Semillero de Astronomía — IEEOH}
}

\setstretch{1.12}
\setlength{\parskip}{0.4em}
\setlength{\parindent}{0pt}
\setcounter{secnumdepth}{3}
\setcounter{tocdepth}{2}
\setlength{\headheight}{26pt}

% -------------------------------------------------
% ESTILOS
% -------------------------------------------------
\titleformat{\chapter}[display]
  {\bfseries\Huge\color{azulCosmos}}
  {\filleft\Large\chaptername\ \thechapter}
  {1ex}
  {\titlerule\vspace{1ex}\filright}
  [\vspace{1ex}\titlerule]

\titleformat{\section}
  {\bfseries\Large\color{azulCosmos}}
  {\thesection}{0.7em}{}

\titleformat{\subsection}
  {\bfseries\large\color{grisTexto}}
  {\thesubsection}{0.6em}{}

\pagestyle{fancy}
\fancyhf{}
\fancyhead[L]{\textcolor{grisSec}{\nouppercase{\leftmark}}}
\fancyhead[R]{\textcolor{grisSec}{Semillero de Astronomía — IEEOH}}
\fancyfoot[C]{\textcolor{grisSec}{\thepage}}
\renewcommand{\headrulewidth}{0.4pt}
\renewcommand{\headrule}{%
  \hbox to\headwidth{\color{azulCosmos}\leaders\hrule height \headrulewidth\hfill}}

\renewcommand{\cftchapfont}{\bfseries\color{azulCosmos}}
\renewcommand{\cftsecfont}{\color{grisTexto}}
\renewcommand{\cftchappagefont}{\bfseries\color{azulCosmos}}
\renewcommand{\cftsecpagefont}{\color{grisTexto}}

% -------------------------------------------------
% ENTORNOS
% -------------------------------------------------
\newtcolorbox{destacado}{
  colback=grisClaro,
  colframe=azulCosmos,
  boxrule=0.8pt,
  arc=2mm,
  left=3mm, right=3mm, top=1.5mm, bottom=1.5mm
}

\newtcbtheorem[number within=chapter]{concepto}{Concepto}{
  enhanced,
  colback=colorConcepto,
  colframe=azulCosmos,
  coltitle=white,
  colbacktitle=azulCosmos,
  attach boxed title to top left={yshift=-2mm, xshift=4mm},
  boxed title style={enhanced, arc=1mm, boxrule=0pt, colframe=azulCosmos},
  fonttitle=\bfseries,
  breakable,
  arc=1.5mm,
  boxrule=0.8pt,
  left=4mm, right=4mm, top=5mm, bottom=2mm,
  separator sign none,
  description delimiters={(}{)}
}{con}

\newtcbtheorem[number within=chapter]{actividad}{Actividad}{
  enhanced,
  colback=colorActividad,
  colframe=azulCosmos,
  coltitle=white,
  colbacktitle=azulCosmos,
  attach boxed title to top left={yshift=-2mm, xshift=4mm},
  boxed title style={enhanced, arc=1mm, boxrule=0pt, colframe=azulCosmos},
  fonttitle=\bfseries,
  breakable,
  arc=1.5mm,
  boxrule=0.8pt,
  left=4mm, right=4mm, top=5mm, bottom=2mm,
  separator sign none,
  description delimiters={(}{)}
}{act}

\newtcbtheorem[number within=chapter]{idea}{Idea}{
  enhanced,
  colback=colorIdea,
  colframe=doradoEstrella,
  coltitle=white,
  colbacktitle=doradoEstrella,
  attach boxed title to top left={yshift=-2mm, xshift=4mm},
  boxed title style={enhanced, arc=1mm, boxrule=0pt, colframe=doradoEstrella},
  fonttitle=\bfseries,
  breakable,
  arc=1.5mm,
  boxrule=0.8pt,
  left=4mm, right=4mm, top=5mm, bottom=2mm,
  separator sign none,
  description delimiters={(}{)}
}{idea}

\newtcbtheorem[number within=chapter]{nota}{Nota}{
  enhanced,
  colback=colorNota,
  colframe=grisSec,
  coltitle=white,
  colbacktitle=grisSec,
  attach boxed title to top left={yshift=-2mm, xshift=4mm},
  boxed title style={enhanced, arc=1mm, boxrule=0pt, colframe=grisSec},
  fonttitle=\bfseries,
  breakable,
  arc=1.5mm,
  boxrule=0.6pt,
  left=4mm, right=4mm, top=5mm, bottom=2mm,
  separator sign none,
  description delimiters={(}{)}
}{nota}

\newtcbtheorem[number within=chapter]{compromiso}{Compromiso}{
  enhanced,
  colback=colorCompromiso,
  colframe=azulCielo,
  coltitle=white,
  colbacktitle=azulCielo,
  attach boxed title to top left={yshift=-2mm, xshift=4mm},
  boxed title style={enhanced, arc=1mm, boxrule=0pt, colframe=azulCielo},
  fonttitle=\bfseries,
  breakable,
  arc=1.5mm,
  boxrule=0.8pt,
  left=4mm, right=4mm, top=5mm, bottom=2mm,
  separator sign none,
  description delimiters={(}{)}
}{comp}

% -------------------------------------------------
% DATOS
% -------------------------------------------------
\newcommand{\NombreSemillero}{Semillero de Astronomía}
\newcommand{\Institucion}{Institución Educativa Enrique Olaya Herrera}
\newcommand{\Ciudad}{Medellín, Colombia}
\newcommand{\AsesorPrincipal}{Daniel Soto Villada}
\newcommand{\CoAsesor}{Gerardo Ríos}
\newcommand{\TituloDocumento}{Clase de Astrobiología\\[0.4em]Ecuación de Drake}

% -------------------------------------------------
% DOCUMENTO
% -------------------------------------------------
\begin{document}
\hypersetup{pageanchor=false}

\begin{titlepage}
    \centering
    \vspace*{1.5cm}
    {\Large\textcolor{grisSec}{\Institucion}\par}
    \vspace{0.2cm}
    {\large\textcolor{grisSec}{\Ciudad}\par}

    \vspace{2.0cm}
    {\Huge\bfseries\textcolor{azulCosmos}{\TituloDocumento}\par}
    \vspace{0.5cm}
    {\Large\textcolor{grisSec}{Semillero de Astronomía}\par}

    \vspace{2.2cm}
    \begin{tcolorbox}[colback=grisClaro,colframe=azulCosmos,boxrule=0.9pt,arc=2mm,width=0.82\textwidth]
        \centering
        {\large\bfseries\textcolor{azulCosmos}{Asesor Principal}\par}
        \vspace{0.2cm}
        {\large \AsesorPrincipal\par}
        \vspace{0.4cm}
        {\large\bfseries\textcolor{azulCosmos}{Co-Asesor}\par}
        \vspace{0.2cm}
        {\large \CoAsesor\par}
    \end{tcolorbox}

    \vfill
    {\large\textcolor{grisSec}{\today}\par}
\end{titlepage}

\frontmatter
\tableofcontents
\clearpage

\mainmatter
\hypersetup{pageanchor=true}

\chapter{Astrobiología y la ecuación de Drake}

\begin{destacado}
    \textbf{Propósito de la clase:} comprender la ecuación de Drake como herramienta de pensamiento científico para discutir la posibilidad de vida inteligente en la galaxia, reconociendo sus límites e implicaciones actuales \cite{baak2026implicationspessimisticlowerlimit,stern2024tectonics}.
\end{destacado}

\section{¿Qué es la astrobiología y por qué usar la ecuación de Drake?}

La astrobiología estudia el origen, evolución, distribución y futuro de la vida en el universo. Esto incluye química prebiótica, habitabilidad planetaria, evolución biológica, geología planetaria, radioastronomía y ciencias de datos. La ecuación de Drake no \emph{demuestra} que exista vida inteligente, pero organiza de forma cuantitativa una gran pregunta: \textit{¿cuántas civilizaciones tecnológicas podrían coexistir hoy en la Vía Láctea?}

\begin{concepto}{Ecuación de Drake como modelo de factores}{con:drake}
La forma clásica de la ecuación es:
\[
N = R_* \, f_p \, n_e \, f_l \, f_i \, f_c \, L
\]
donde:
\begin{itemize}[leftmargin=1.5em]
    \item $R_*$: tasa de formación de estrellas apropiadas en la galaxia.
    \item $f_p$: fracción de esas estrellas con sistemas planetarios.
    \item $n_e$: número medio de planetas potencialmente habitables por sistema.
    \item $f_l$: fracción de planetas habitables donde aparece vida.
    \item $f_i$: fracción de planetas con vida donde surge inteligencia.
    \item $f_c$: fracción de civilizaciones inteligentes que desarrollan comunicación detectable.
    \item $L$: tiempo promedio (años) durante el cual una civilización emite señales detectables.
\end{itemize}
\end{concepto}

\section{Lectura moderna: optimismo, pesimismo y paradoja de Fermi}

El punto crítico en la discusión contemporánea es la enorme incertidumbre de los factores biológicos y sociotecnológicos. Estudios recientes resaltan que, incluso con supuestos astronómicos favorables, basta que algunos factores sean muy pequeños para que $N$ se reduzca drásticamente \cite{baak2026implicationspessimisticlowerlimit}.

\begin{nota}{Relación con la paradoja de Fermi}{nota:fermi}
Si muchas civilizaciones deberían existir, ¿por qué no las observamos? Esta tensión entre expectativa teórica y ausencia de evidencia observacional se conoce como paradoja de Fermi. Una posible respuesta es que algunos factores de la ecuación (por ejemplo $f_i$, $f_c$ o $L$) pueden ser extremadamente bajos.
\end{nota}

Además, se propone ampliar la ecuación incorporando restricciones geofísicas para vida compleja, como la coexistencia de continentes y océanos, y tectónica de placas sostenida en escalas geológicas \cite{stern2024tectonics}. Bajo esa visión, no basta con estar en zona habitable: también importan condiciones planetarias de larga duración.

\section{Interpretación detallada de cada término}

\subsection{$R_*$: formación estelar}
Este factor ha mejorado gracias a observaciones galácticas. Suele tratarse como relativamente mejor conocido que los factores biológicos. No obstante, no todas las estrellas son igual de favorables para estabilidad climática planetaria.

\subsection{$f_p$ y $n_e$: arquitectura planetaria}
Con descubrimientos de exoplanetas, hoy se sabe que los sistemas planetarios son comunes. El reto no es encontrar planetas, sino caracterizar cuáles mantienen agua líquida, atmósferas y estabilidad orbital suficiente para procesos evolutivos largos.

\subsection{$f_l$: aparición de vida}
Es uno de los términos más inciertos. La Tierra aporta un único ejemplo confirmado. No sabemos si la vida emerge con alta probabilidad cuando las condiciones químicas son adecuadas, o si es un evento rarísimo.

\subsection{$f_i$: transición a inteligencia}
La vida microbiana pudo surgir relativamente temprano en la Tierra, pero la inteligencia tecnológica tardó miles de millones de años. Esto sugiere que el paso de vida simple a vida compleja puede ser un cuello de botella.

\subsection{$f_c$: tecnoseñales detectables}
No toda inteligencia desarrolla radioastronomía, telecomunicaciones o emisiones detectables desde lejos. También es posible que usen tecnologías que no estamos buscando.

\subsection{$L$: duración de detectabilidad}
Este término domina muchos escenarios: aunque surjan civilizaciones, si su fase detectable es corta, la superposición temporal entre ellas puede ser mínima.

\section{Ejemplos desarrollados}

\begin{actividad}{Ejemplo 1: escenario moderadamente optimista}{act:ejemplo1}
Supongamos:
\[
R_* = 2,\ f_p = 0.8,\ n_e = 0.3,\ f_l = 0.3,\ f_i = 0.02,\ f_c = 0.2,\ L = 8{,}000
\]
Entonces:
\[
N = 2 \cdot 0.8 \cdot 0.3 \cdot 0.3 \cdot 0.02 \cdot 0.2 \cdot 8000 = 4.608
\]
Interpretación: podrían coexistir alrededor de 5 civilizaciones detectables en la Vía Láctea.
\end{actividad}

\begin{actividad}{Ejemplo 2: escenario pesimista}{act:ejemplo2}
Si mantenemos los términos astronómicos pero reducimos factores biológicos y de longevidad:
\[
R_* = 2,\ f_p = 0.8,\ n_e = 0.3,\ f_l = 0.01,\ f_i = 0.001,\ f_c = 0.1,\ L = 500
\]
resulta:
\[
N = 2 \cdot 0.8 \cdot 0.3 \cdot 0.01 \cdot 0.001 \cdot 0.1 \cdot 500 = 0.00024
\]
Interpretación: la expectativa de civilizaciones coexistentes es mucho menor que 1; en términos prácticos, podría no haber ninguna detectable en este momento \cite{baak2026implicationspessimisticlowerlimit}.
\end{actividad}

\begin{actividad}{Ejemplo 3: extensión geofísica inspirada en literatura reciente}{act:ejemplo3}
Tomamos el escenario del ejemplo 1 e incorporamos dos factores adicionales:
\[
N' = N \cdot f_{oc} \cdot f_{pt}
\]
con $f_{oc}$ = fracción de planetas habitables con continentes y océanos significativos, y $f_{pt}$ = fracción de esos planetas con tectónica de placas activa por al menos 0.5 Ga \cite{stern2024tectonics}.

Si $f_{oc} = 0.1$ y $f_{pt} = 0.01$:
\[
N' = 4.608 \cdot 0.1 \cdot 0.01 = 0.004608
\]
Interpretación: incluir filtros geodinámicos puede reducir fuertemente la estimación de civilizaciones complejas.
\end{actividad}

\section{Discusión didáctica para el estudiante}

\begin{idea}{Idea central de la clase}{idea:idea-central}
La ecuación de Drake no es una calculadora exacta del universo; es un mapa de ignorancias y evidencias. Cada término indica qué sabemos bien, qué conocemos a medias y qué apenas empezamos a investigar.
\end{idea}

\begin{nota}{Error común}{nota:error}
Un valor alto o bajo de $N$ no es ``la verdad final''. Es un resultado condicionado por supuestos. Cambiar supuestos razonables puede alterar órdenes de magnitud completos.
\end{nota}

\section{Preguntas guía para debate en clase}

\begin{itemize}[leftmargin=1.5em]
    \item ¿Qué término de la ecuación consideras más incierto y por qué?
    \item Si encontráramos vida microbiana fuera de la Tierra, ¿qué término cambiaría primero?
    \item ¿Es posible que existan civilizaciones no detectables por nuestros métodos actuales?
    \item ¿Cómo se relaciona la sostenibilidad de una civilización con el parámetro $L$?
    \item ¿Crees que los factores geológicos (tectónica, continentes, océanos) deben incorporarse siempre en modelos de vida compleja?
\end{itemize}

\section{Tareas para el estudiante}

\begin{compromiso}{Trabajo autónomo}{comp:tareas}
\begin{enumerate}[leftmargin=1.5em]
    \item \textbf{Cálculo base:} proponga un conjunto propio de valores para los 7 factores clásicos y calcule $N$ paso a paso.
    \item \textbf{Análisis de sensibilidad:} cambie solo uno de los términos ($f_l$, $f_i$ o $L$) en un factor de 10 y explique cuánto cambia $N$.
    \item \textbf{Comparación de modelos:} use su resultado base y luego incorpore $f_{oc}$ y $f_{pt}$ para calcular $N'$. Escriba una comparación crítica entre ambos modelos.
    \item \textbf{Miniensayo (1–2 páginas):} responda: \enquote{¿La paradoja de Fermi sugiere que estamos solos o que todavía buscamos mal?}. Debe citar al menos una idea de \cite{baak2026implicationspessimisticlowerlimit} y una de \cite{stern2024tectonics}.
    \item \textbf{Pregunta abierta de investigación:} formule una pregunta que pueda investigarse en el semillero con recursos escolares (por ejemplo, tecnoseñales, exoplanetas habitables, o evolución de vida compleja).
\end{enumerate}
\end{compromiso}

\backmatter
\chapter*{Referencias}
\addcontentsline{toc}{chapter}{Referencias}

\begin{thebibliography}{99}

\bibitem{baak2026implicationspessimisticlowerlimit}
Baak, M., y Snoek, H. (2026).
\textit{Implications of the Pessimistic Lower Limit on the Drake Equation}.
arXiv:2603.02372.
\url{https://arxiv.org/abs/2603.02372}

\bibitem{stern2024tectonics}
Stern, R. J., y Gerya, T. V. (2024).
The importance of continents, oceans and plate tectonics for the evolution of complex life:
implications for finding extraterrestrial civilizations.
\textit{Scientific Reports}, 14, 8552.
\url{https://doi.org/10.1038/s41598-024-54700-x}

\end{thebibliography}

\end{document}
